\subsection{对初等自同构的应用}

命题 9. 设 k 是一个特征为 0 的域，g 是 k 上的一个李代数。若 \(a \in \operatorname{Aut}_{e}(\mathfrak{g})\)，则 a - 1 的幂零空间的维数大于或等于 g 的秩。

由“Lefschetz 原理”（《代数》，第五章，§14，第 6 节，定理 5 的推论 2），k 是以 C 为扩张域的一族子域 \((k_{i})_{i\in\mathrm{I}}\) 的上升定向并。设 \((e_{\alpha})\) 是 g 在 k 上的一个基，\(x_{1},\ldots,x_{m}\) 是 g 的元素，使得 \(\operatorname{ad}(x_{1}),\ldots,\operatorname{ad}(x_{m})\) 幂零且 \(a=e^{\operatorname{ad}(x_{1})}\ldots e^{\operatorname{ad}(x_{m})}\)。设 \(c_{\alpha\beta}^{\gamma}\) 是 g 关于基 \((e_{\alpha})\) 的结构常数，并设 \((x_{r}^{\alpha})\) 是 \(x_{r}\) 关于这个基的分量 \((1\leq r\leq m)\)。存在一个指标 \(j\in I\)，使得诸 \(c_{\alpha\beta}^{\gamma}\) 与诸 \(x_{r}^{\alpha}\) 都属于 \(k_{j}\)。设 \(g_{j}=\sum_{\alpha}k_{j}e_{\alpha}\)；这是 \(k_{j}\) 上的一个李代数，包含 \(x_{1},\ldots,x_{m}\)，且 \(a_{j}\)（a 在 \(g_{j}\) 上的限制）是 \(g_{j}\) 的一个初等自同构。\(a_{j}\) 到 \(g_{j}\otimes_{k_{j}}C\) 上的扩张是 \(a_{j}\otimes1\)（\(g_{j}\otimes C\) 的一个初等自同构）。于是，设 \(G_{j}\) 是一个以 \(g_{j}\otimes C\) 为李代数的连通复李群，s 是 \(G_{j}\) 的一个元素，使得 \(\operatorname{Ad}(s)=a_{j}\otimes1\)。把命题 8 应用于二元组 \((G_{j},s)\)，就得到 \(a_{j}\otimes1-1\) 的幂零空间是 n 维的，于是

\[
n \geq \operatorname{rk} (\mathfrak {g} _ {j} \otimes \mathbf {C}) = \operatorname{rk} (\mathfrak {g} _ {j}) = \operatorname{rk} (\mathfrak {g}).
\]

但这个幂零空间与 \(a_{j} - 1\) 的以及 \(a - 1\) 的幂零空间有相同的维数。命题得证。

